\section{Finite-deformation coefficient from material conservation}

The computational reference configuration is that of the solid skeleton.  Let
\(\mbf F_s\) denote the skeleton deformation gradient and
\(J_s=\det\mbf F_s\).  For the registered load-bearing solid phases
\(a\in\mc S\), let \(\phi_a\) be the current phase volume fraction and
\(\rho_a^\alpha\) the current bulk density of component \(\alpha\) in phase
\(a\).  Summing the solid phase/component material-conservation storage terms
gives the intrinsic skeleton specific volume
\begin{equation}
 \bar v_s=
 \frac{J_s\displaystyle\sum_{a\in\mc S}\phi_a}
 {\displaystyle\sum_{a\in\mc S}\sum_\alpha J_s\rho_a^\alpha}.
 \label{eq:summed-specific-volume}
\end{equation}
The denominator is a referential component accumulation.  Dividing it by
\(J_s\) gives a bulk solid partial density; the intrinsic skeleton density is
instead
\begin{equation}
 \bar\rho_s=\frac{1}{\bar v_s}=
 \frac{\displaystyle\sum_{a\in\mc S}\sum_\alpha J_s\rho_a^\alpha}
 {J_s\displaystyle\sum_{a\in\mc S}\phi_a}.
 \label{eq:intrinsic-density}
\end{equation}

The nonlinear coefficient is
\begin{equation}
 B=1-\frac{1}{v_{s0}}
 \left.\frac{\partial\bar v_s}{\partial J_s}\right|_{\bar p_E,\theta_{\mc S},\ldots},
 \label{eq:nonlinear-biot}
\end{equation}
where \(v_{s0}\) is the intrinsic skeleton reference specific volume and
\(\bar p_E\) is equivalent pore pressure. The qualifier ``fixed
\(\bar p_E\)'' is a
constitutive partial derivative.  It is not the total derivative obtained by
allowing every finite-element field to vary along an arbitrary solution path.

The reference stress consumed by the mechanics residual is
\begin{equation}
 \mbf P_s=\mbf P_s^{\prime\prime}
 -B \bar p_E J_s\mbf F_s^{-T},
 \label{eq:reference-biot-stress}
\end{equation}
where \(\mbf P_s^{\prime\prime}\) is the effective first
Piola--Kirchhoff stress evaluated at fixed \(\bar p_E\). Because \(B\) is a state
variable, the consistent Newton derivative of \cref{eq:reference-biot-stress}
contains derivatives of \(B\), \(\bar p_E\), \(J_s\), and \(\mbf F_s^{-T}\).

In the single-solid infinitesimal-strain limit, evaluation at the reference
state gives
\begin{equation}
 B_0=1-\frac{K}{K_s}.
 \label{eq:classical-limit}
\end{equation}
The implementation is required to approach \cref{eq:classical-limit} under
shrinking perturbations rather than merely reproduce it at one selected state.
